\section{Introduction}\label{sec:intro}

A scheme of rank $r$ for $n\times m$ by $m\times p$ matrix multiplication computes the product with $r$
non-commutative multiplications of linear forms; equivalently it writes the matrix-multiplication
tensor $\MM nmp$ as a sum of $r$ rank-one tensors \cite{burgisser1997algebraic,landsberg2012tensors,blaser2013fast}.
Since Strassen's rank-7 scheme for $\MM222$ \cite{strassen1969gaussian}, upper bounds for small formats have come from
clever constructions \cite{hopcroft1971minimizing,laderman1976noncommutative,smirnov2013bilinear}, from
SAT and local search \cite{heule2019local,heule2021new}, from reinforcement learning \cite{fawzi2022discovering}, from
program evolution \cite{novikov2025alphaevolve} and, most productively in recent years, from random walks on \emph{flip graphs}
\cite{kauers2023flip,arai2024adaptive,moosbauer2025flip,kauers2025meta,perminov2025ternary}.
The field matters: over $\F$ there are schemes that do not exist in other characteristics.
The best-known example is the rank-47 scheme for $\MM444$ \cite{fawzi2022discovering}; for $\MM555$ the $\F$ record
went $98\to96\to95\to94\to93$ \cite{sedoglavic2021tensor,fawzi2022discovering,kauers2022still,arai2024adaptive,moosbauer2025flip}.

Best-known ranks for all formats up to $32\times32\times32$ are collected in two public tables
\cite{sedoglavic2026catalogue,lacelle2026catalog}; two repositories hold explicit schemes for the formats up to $16\times16\times16$
\cite{perminov2026fmmrepo} and for 25 formats \cite{witteveen2026schemes}.
Small formats are the subject of intense search.
Large formats mostly are not searched at all: the listed ranks there come from composing smaller
schemes by block sums and Kronecker products, so they respond to any improvement of a component,
and the components that matter in characteristic two (the rank-47 scheme above, in particular) are
not the ones the general-ring tables were built from.

This paper describes a pipeline that attacks that regime and the engineering needed to make its
output believable.
It was used for a campaign in September 2026. The results, pooled from three stages of the project (July, 2026-09-22 to 09-26,
and the campaign of 2026-09-26 to 09-29), are 795 shapes for which our best verified scheme is below every GF(2)-valid rank listed in four public tables at the pinned revisions
(733 of them also below every rank listed in any ring); they are reported in a companion paper \cite{peterson2026results}.
Here we concentrate on \emph{how} the schemes were obtained, how they were checked, and what each stage of the campaign was worth.
The campaign ran 4{,}816 walks totalling \(3.1\times10^{11}\) proposals over 2{,}708 distinct shapes,
materialized 942 composition products, replayed 1{,}352 derivation recipes from the matmulcatalog (1{,}098 of them exactly over $\F$),
and verified every tensor that is claimed.

\paragraph{Contributions.}
\begin{enumerate}
\item A \textbf{directed flip-graph walker} for formats with all sides at least 2 and factors of up to $1{,}024$ bits
(\cref{sec:walker}): proposals are sampled only among factor buckets that contain a flippable pair,
factors are multiword bit vectors, memory is allocated once, and a walk is a deterministic function of
(seed file, number of proposals, nonce). We report its per-proposal cost as a function of factor width
and the byte-for-byte replay of 42 distinct campaign walks.
\item A \textbf{cross-shape derivation step} (\cref{sec:derive}): coordinate deletion applied to the best
scheme after a shared-factor basis rewrite, followed by shared-factor compression. Walks from these
``projection children'' lowered the rank below the child's own starting rank in 792 of 880 cases, at 44M proposals per
lowering against 0.3--9G for the seed classes that do not descend from projection children (0.16G for the deepening of projection lineages); children start unreduced, so this measures the slack that projection
leaves for the walker rather than the yield of records.
\item A \textbf{rank-only composition closure} with witness-backed materialization (\cref{sec:composition}):
a dynamic programme over all 5{,}984 sorted formats with sides $\le32$ that prices block sums and Kronecker products in
0.3\,s, followed by construction and exact checking of each priced improvement.
\item \textbf{Seeds from public tables} (\cref{sec:seeds}): reduction of published ternary and integer schemes
modulo two behind a ring gate, and replay of matmulcatalog derivation recipes over $\F$, including detection of recipes whose
``F2'' stamp over-claims and a case study of a cache-history dependence that made replays
non-reproducible until it was fixed.
\item A \textbf{scoring and claim discipline} (\cref{sec:scoring}): what to compare against, which comparisons are
invalid (commutative algorithms, wrong ring, stale tables), and how field-agnostic tables distort the picture.
\item \textbf{Lineage bundles and a verification protocol} (\cref{sec:lineage}) in which every tensor is expanded against
the full matrix-multiplication tensor by two separately written implementations; projection, public-seed, archived and retained rows
are recomputed from their source; 42 of 244 distinct walks are replayed and 71 of 82 catalogue replays are reproduced from an empty cache;
and the check is repeated from a clean extraction.
\item An \textbf{empirical account of yield and cost} per stage, including the negative results (\cref{sec:eval}).
\end{enumerate}

\begin{figure}[t]
\centering
\resizebox{\textwidth}{!}{%
\begin{tikzpicture}[
  box/.style={draw,rounded corners=2pt,align=center,inner sep=4pt,font=\small,minimum height=10mm,minimum width=46mm},
  shade/.style={fill=gray!14},
  arr/.style={-{Stealth[length=2.2mm]},thick}]
  % sources
  \node[box] (tab) at (0,6.0) {four public tables\\ {\footnotesize Lille, matmulcatalog, FMM, fmm-schemes}};
  \node[box] (pub) at (6.0,6.0) {integer/ternary schemes\\ reduced mod 2 {\footnotesize(\S\ref{sec:seeds})}};
  \node[box] (rep) at (12.0,6.0) {matmulcatalog recipes\\ replayed over $\mathbb F_2$ {\footnotesize(\S\ref{sec:seeds})}};
  % scoreboard
  \node[box,shade,minimum width=76mm] (board) at (6.0,3.7) {scoreboard: best verified scheme per sorted format {\footnotesize(\S\ref{sec:scoring})}};
  % operators
  \node[box] (clo) at (0,1.4) {composition closure\\ price $\to$ materialize {\footnotesize(\S\ref{sec:composition})}};
  \node[box] (proj) at (6.0,1.4) {projection cascade\\ rewrite, project, compress {\footnotesize(\S\ref{sec:derive})}};
  \node[box] (walk) at (12.0,1.4) {directed flip walks\\ $4\!\times\!10^{7}$--$2\!\times\!10^{8}$ proposals {\footnotesize(\S\ref{sec:walker})}};
  % gate and bundle
  \node[box] (gate) at (3.0,-1.2) {exactness gate\\ Python; Ruby audit {\footnotesize(\S\ref{sec:repr})}};
  \node[box,shade] (bundle) at (9.0,-1.2) {lineage bundles; clean-extraction check\\ {\footnotesize(\S\ref{sec:lineage})}};
  \draw[arr] (tab) -- node[left,font=\footnotesize,pos=0.55]{bars} (board.north west);
  \draw[arr] (pub) -- (board.north);
  \draw[arr] (rep) -- (board.north east);
  \draw[arr] (board.south west) -- node[left,font=\footnotesize]{leads} (clo.north);
  \draw[arr] (board.south) -- node[right,font=\footnotesize]{parents} (proj.north);
  \draw[arr] (board.south east) -- node[right,font=\footnotesize]{seeds} (walk.north);
  \draw[arr] (clo.south) -- (gate.north west);
  \draw[arr] (proj.south) -- (gate.north east);
  \draw[arr] (walk.south) |- ++(0,-0.75) -| (gate.north east);
  \draw[arr] (gate) -- (bundle);
  \draw[arr,dashed] (gate.west) -- ++(-3.9,0) |- (board.west);
\end{tikzpicture}}
\caption{The pipeline. Every tensor enters the scoreboard only through the exactness gate (dashed arrow); the public tables steer effort (bars, leads, parents, seeds) but never admission.}
\label{fig:pipeline}
\end{figure}


\paragraph{What is not claimed.}
Projection, block composition and Kronecker products are classical; our contribution there is the
shared-factor basis rewrite before projection, the materialization from GF(2)-valid leaves and the measured yield, not a new mathematical operation.
Every rank reported anywhere is an upper bound. We prove no lower bounds, we do not claim optimality, and, following
the record policy of the project, results that beat a table are called \emph{apparent} until a complete prior-art
audit has been done (\cref{sec:scoring,sec:limits}).

\paragraph{Organization.}
\Cref{sec:background} fixes notation. \Cref{sec:repr} defines the scheme format and the exactness gate,
\cref{sec:walker} the walker, \cref{sec:derive} projection and basis rewriting, \cref{sec:composition} composition closure and
\cref{sec:seeds} the use of public schemes. \Cref{sec:scoring} covers scoring, \cref{sec:lineage} provenance and verification,
\cref{sec:eval} the campaign's measured behaviour and \cref{sec:limits} limitations.
